\chapter{Preface}

Throughout this textbook, \textbf{logical formulas} are interpreted according to the following conventions:
\begin{itemize}
    \item \textbf{Parentheses} have the highest precedence. For example,
    \[
    (P \lor Q) \land R
    \]
    is evaluated by first computing $P \lor Q$.
    \item The \textbf{logical connectives} are interpreted in the order
    \[
    \neg \; > \; \land \; > \; \lor \; > \; \implies \; > \; \iff
    \]
    Hence
    \[
    \neg P \land Q \lor R \implies S \iff T
    \]
    is understood as
    \[
    ((((\neg P) \land Q) \lor R) \implies S) \iff T
    \]
    \item A \textbf{quantifier} binds the entire formula that follows it. Thus
    \[
    \forall x, P(x) \land Q(x) \implies R(x)
    \]
    means
    \[
    \forall x\, ((P(x)\land Q(x))\implies R(x))
    \]
    \item The symbol $\Downarrow$ is \textbf{not} a logical connective;
    it is used only to indicate the flow of a proof or an explanation. For example,
    \[
    \begin{array}{c}
        P \\
        \Downarrow \\
        Q
    \end{array}
    \]
    merely expresses that $P$ is true, and therefore $Q$ is true in the current discussion;
    it is not itself a proposition.
\end{itemize}
